\documentclass[output=paper]{langscibook}
\ChapterDOI{10.5281/zenodo.22880268}
\title{Introduction}
\author{Ulrike Mosel}
\abstract{}
\begin{document}
\maketitle\label{chap:21}
\begin{figure}[H]
    \includegraphics[trim=0 0 0 0,clip,width=.9\linewidth]{figures/21_1.jpg}
    \caption{\textit{Hua gunaha} – Paddling ashore}
    \label{fig:21_1}
\end{figure}
Apart from one modern kind of metal boat, I only observed dug-out canoes with outriggers in the Teop area during my fieldwork trips between 2000 and 2014. There are two kinds of dug-out outrigger canoes: the indigenous outrigger canoe called \textit{sinivi} \textit{suuna} and the \textit{potee}, a new type of outrigger canoe that is said to originate from the Solomon Islands. Both kinds of canoe are dugouts built from a single log, but differ in the width of the hull and the construction of the outrigger.
\begin{figure}[H]
    \includegraphics[trim=0 0 0 0,clip,width=.9\linewidth]{figures/21_2.jpg}
    \caption{\textit{Sinivi suuna} -- Indigenous canoe}
    \label{fig:21_2}
\end{figure}
\begin{figure}[H]
    \includegraphics[trim=0 0 0 0,clip,width=.9\linewidth]{figures/21_3.jpg}
    \caption{\textit{Potee} -- A new type of canoe}
    \label{fig:21_3}
\end{figure}
The first text by Mark Tahi and Enoch Horai Magum describes the hull of the \textit{sinivi} \textit{suuna} and is followed by Ruth Sima Rigamu’s description of the outrigger float, which is an excerpt from her texts about trees. The tree in question is called \textit{kopihoo}. As we could not identify its scientific or its English name, we included two photographs of this tree. In third chapter Paul Morekevan and Enoch Horai Magum describe the hull and the outrigger of the \textit{potee} canoe. The word \textit{sinivi} is a generic term for both types of canoe, and the name for the Blackboard Tree, see Chapter \ref{chap:7}, which is used for building the hull of these canoes. 

The construction of the Teop \textit{sinivi} \textit{suuna}  is similar to the outrigger dug-out canoes on Petats, an island off the west coast of Buka, and in Kurtatchi, a village on the north coast of Bougainville that have been described by \textcite[360-381]{Blackwood:1935im}. Her list of ``terms used in connection with dug-out canoe'' shows several terms of the Kurtachi language that are closely related to the corresponding Teop terms, see \tabref{tab:canoe} and \textcite[367]{Blackwood:1935im}.
\begin{table}[h]
    \caption{Similar Kurtatchi and Teop terms related to the dugout outrigger canoe}
    \label{tab:canoe}
    \begin{tabularx}{.9\textwidth}{lll}
        \lsptoprule
            \textsc{meaning of term}&\textsc{kurtatchi}&\textsc{teop}\\
        \midrule
            Dug-out canoe with outrigger&\textit{tʃinív}&\textit{sinivi}\\
            Tree of which it is generally  made&\textit{tʃinív}&\textit{sinivi}\\
            Outrigger booms and pole forming  platform&\textit{wanér}&\textit{vanai}\\
            Outrigger float&\textit{ag} or \textit{waga}&\textit{aga}\\
        \lspbottomrule
    \end{tabularx}
\end{table}
According to the people \textcite[360]{Blackwood:1935im} and \textcite[186]{Shoffner:1976eg} worked with, the outrigger dugout canoe was introduced to North-Bougainville before European contact. In contrast to other people of Bougainville, the Teop people did not build plank canoes (\cite[370]{Blackwood:1935im}, \cite[186]{Shoffner:1976eg}).

The \textit{potee} canoe, see \figref{fig:21_3}, is not regarded as a \textit{sinivi} \textit{suuna} `indigenous canoe', but said to be a type of canoe from the Solomon Islands to which the Teop people attached an outrigger with two booms, see Chapter \ref{chap:24}:\ref{ex:text18-47}ff. The terms used for the parts of the \textit{potee} and the process of its building are the same as for the \textit{sinivi} \textit{suuna}. When I stayed at Hiovabon for my last fieldwork in 2014, most people used the \textit{sinivi} \textit{suuna}.
\begin{figure}[H]
    \includegraphics[trim=0 0 0 0,clip,width=.9\linewidth]{figures/21_4.jpg}
    \caption{Teop Harbour in 2013}
    \label{fig:21_4}
\end{figure}
The two texts about the \textit{sinivi} \textit{suuna} and the \textit{potee} do not mention which kinds of wood are used to make the hull and the outriggers, but Joyce Maion describes in Chapter \ref{chap:7} the \textit{sinivi} `Blackboard Tree (\textit{Alstonia} \textit{scholaris})' and mentions two further trees: the \textit{maramakasi} which we did not identify and the \textit{bio} a kind of Ficus tree.
Other trees for canoe building are:
\begin{itemize}
    \item \textit{Marakaoto} `tree species (\textit{Terminalia} \textit{species})' -- \url{https://dictionaria.clld.org/units/teop-marakaoto}
    \item \textit{Tako} `Alexandrian Laurel (\textit{Callophyllum} \textit{mophyllum})' -- \url{https://dictionaria.clld.org/units/teop-tako_1}
    \item \textit{Voata} `tree species (\textit{Euodia} \textit{elleryana})' -- \url{https://dictionaria.clld.org/units/teop-voata}
\end{itemize}
\end{document}